% Chapter 1 draft: General biological introduction.
% Internal provenance and chronology audit:
%   research_audit/ch01_docs.md
% Current implementation anchors:
%   DOC0321
%   main/bioseq/backend/sequences/sequences.cppm
%   main/bioseq/io/{fasta,fastq,pipeline,public_api}.cppm
%   main/bioseq/sql/bioseq--1.0.sql
% Scholarly references are given as full footnotes so this chapter remains
% readable independently. Primary sources are also registered in the thesis
% bibliography where they are reused by later chapters.

\section{Biological sequences as information-bearing observations}
\label{sec:biology-sequences}

The data processed by BioSeqDB originate in molecules, but an alignment
algorithm does not receive a molecule. It receives an ordered sequence of
symbols. This distinction is elementary and important. A physical DNA molecule
has chemical structure, direction, complementary strands and a history inside
a cell. A database value is a finite string together with metadata that
describes how the string was obtained and how its symbols should be
interpreted. A reliable computational system must preserve enough of that
interpretation to distinguish, for example, a nucleotide sequence from a
protein sequence and an unknown base from an ordinary base.

\subsection{DNA, RNA and protein alphabets}

DNA is a polymer whose monomers contain one of four canonical bases:
adenine, cytosine, guanine and thymine. Their conventional one-letter symbols
form the alphabet
\[
  \Sigma_{\mathrm{DNA4}} =
  \{\mathtt{A},\mathtt{C},\mathtt{G},\mathtt{T}\}.
\]
The order of symbols is normally written from the \(5'\) end of a strand to its
\(3'\) end. DNA in a chromosome is usually double stranded. Adenine pairs with
thymine, while cytosine pairs with guanine. Consequently, the sequence
\(\mathtt{ACGTT}\) on one strand determines
\(\mathtt{AACGT}\) on the opposite strand when that strand is also written in
the \(5'\)-to-\(3'\) direction. The second string is the \emph{reverse
complement} of the first. A mapper cannot assume that an observed read has the
same orientation as the stored reference. It commonly has to consider both the
forward sequence and its reverse complement.

Four symbols are sufficient for an unambiguous DNA string, but experimental
data are not always unambiguous. The symbol \(\mathtt{N}\) conventionally means
that the nucleotide at a position is unknown or was not resolved. Other
International Union of Pure and Applied Chemistry (IUPAC) codes describe sets
of possible bases, such as a purine at a position. These
are uncertainty symbols, not additional chemical bases. Treating
\(\mathtt{N}\) as if it were an ordinary fifth base would change the meaning of
exact matches, hashes and alignment scores. A sequence type therefore needs an
explicit alphabet and an explicit policy for symbols outside that alphabet.
The current BioSeqDB sequence core follows this principle at the software level:
its compact ranges are parameterized by an alphabet, and current code paths
instantiate both four-symbol and five-symbol DNA representations. The later
sequence-abstraction chapter explains how those symbols are packed and viewed.

RNA uses adenine, cytosine, guanine and uracil, so a basic RNA alphabet replaces
\(\mathtt{T}\) with \(\mathtt{U}\). Many RNA molecules are transcribed from
DNA. Protein-coding transcripts can then be translated in groups of three
nucleotides, called codons, into amino acids. For protein-coding genes,
transcription followed by translation yields the familiar DNA-to-RNA-to-protein
flow. Crick's central dogma is the narrower restriction that residue-by-residue
sequence information cannot flow from a protein sequence into either another
protein sequence or a nucleic acid.\footnote{F. Crick, ``Central dogma of
molecular biology'',
\emph{Nature} 227 (1970), pp.~561--563,
\url{https://doi.org/10.1038/227561a0}.} It is not a claim that all RNA is
translated. Ribosomal, transfer, regulatory and many other RNAs are functional
without becoming proteins.

A protein sequence is written over the twenty common amino-acid symbols.
Protein comparison differs from nucleotide comparison in two ways that matter
computationally. First, the alphabet is much larger. Second, substitutions are
not naturally divided into only ``same'' and ``different''. Replacing an amino
acid with one having similar chemical properties may be more plausible than a
replacement with a very different residue. Protein aligners therefore often
use empirical substitution matrices. Nucleotide aligners can also use a
matrix, but a simpler match, mismatch and gap model is often appropriate.
Table~\ref{tab:biological-alphabets} summarizes the distinctions that storage
and alignment must preserve.

\begin{table}[ht]
  \centering
  \small
  \begin{tabular}{p{0.16\linewidth}p{0.27\linewidth}p{0.47\linewidth}}
    \toprule
    Sequence class & Core symbols & Computational consequences \\
    \midrule
    DNA &
    \(\mathtt{A,C,G,T}\), with optional ambiguity codes &
    Reverse-complement orientation is biologically meaningful. Unknown and
    degenerate symbols need explicit matching and hashing policies. \\
    RNA &
    \(\mathtt{A,C,G,U}\), with optional ambiguity codes &
    Transcript sequences may cross splice junctions. A system must state
    whether it preserves \(\mathtt{U}\) or normalizes it to
    \(\mathtt{T}\). \\
    Protein &
    Twenty common amino-acid symbols, plus data-format conventions for
    ambiguity and exceptional residues &
    A larger alphabet requires more bits per symbol. Substitution scores
    normally encode biochemical or evolutionary similarity rather than
    equality alone. \\
    \bottomrule
  \end{tabular}
  \caption{Biological alphabets and the distinctions that a sequence-processing
    system must not erase. The table describes biological data in general. It
    does not assert that every alphabet or ambiguity convention is exposed by
    every BioSeqDB SQL path.}
  \label{tab:biological-alphabets}
\end{table}

The alphabet is thus part of the type of a sequence, not a decorative label.
It determines the set of valid symbols, the minimum number of bits needed for
storage, whether a complement operation exists and how comparison should
behave. This observation connects basic biology to a central systems decision
in BioSeqDB: algorithms should consume a common sequence abstraction while
remaining explicit about the alphabet represented by that abstraction.

\section{From genomes to computational records}
\label{sec:genomes-to-records}

\subsection{Genomes, chromosomes and assemblies}

A \emph{genome} is the genetic material carried by an organism or biological
sample. Depending on the organism and the question, it may be distributed
among one or more chromosomes, plasmids or organellar molecules. The genome of
a diploid individual contains homologous chromosome copies that need not be
identical. A microbial sample can contain many species and strains at once.
Consequently, ``the genome'' is not always one linear string.

Sequencing and assembly software nevertheless have to produce finite records.
A \emph{contig} is a contiguous consensus sequence assembled from overlapping
evidence. Scaffolds can order or orient contigs while retaining unresolved
gaps. A \emph{reference assembly} is a selected set of such sequences used as a
coordinate system and comparison target. It is extremely useful, but it is not
the unique biological truth for every sample. Regions absent from the
reference, alternative haplotypes and repeated segments can all make a
reference-centered explanation incomplete.

This distinction affects the interpretation of an alignment. Mapping a read
to a reference says that the chosen reference interval is a plausible source
or homologous locus under a scoring model. It does not prove that the sample
is identical to the reference at that interval. It also does not prove that
the reported location is unique when multiple loci are sufficiently similar.
The alignment result must therefore retain coordinates, orientation, score
and, when available, information about competing placements.

\subsection{Reads are measurements, not perfect substrings}

A sequencing instrument observes molecules and emits \emph{base calls}. A
base call is the instrument's interpretation of a physical signal. The calls
for one observed molecule or molecule fragment form a \emph{read}. Some
protocols produce relatively short reads, others produce long reads, and
paired-end protocols retain a relationship between two reads from opposite
ends of a fragment. The categories are operational rather than absolute. What
matters to an aligner is the joint profile of read length, error, orientation
and expected biological divergence.

FASTA stores named sequence records. FASTQ additionally associates each called
symbol with a quality character. The widely used Phred convention expresses
an estimated error probability \(p\) as
\[
  Q = -10\log_{10}p.
\]
The FASTQ format and the historical variation in its quality encodings are
documented by Cock et al.\footnote{P.~J.~A. Cock, C.~J. Fields, N.~Goto,
M.~L. Heuer and P.~M. Rice, ``The Sanger FASTQ file format for sequences with
quality scores, and the Solexa/Illumina FASTQ variants'',
\emph{Nucleic Acids Research} 38(6), 2010, pp.~1767--1771,
\url{https://doi.org/10.1093/nar/gkp1137}.} A quality score is useful evidence,
but it is neither a guarantee nor a biological property of the base itself.
Calibration can vary by instrument, chemistry and processing pipeline.

Sequencing errors are commonly described as substitutions, insertions and
deletions. A substitution reports one symbol in place of another. An insertion
adds a symbol relative to the underlying molecule or comparison target. A
deletion omits one. Long-read and short-read technologies can have different
mixtures and contexts of these errors, which is why one universal scoring
preset is insufficient for all data.\footnote{S. Goodwin, J.~D. McPherson and
W.~R. McCombie, ``Coming of age: ten years of next-generation sequencing
technologies'', \emph{Nature Reviews Genetics} 17 (2016), pp.~333--351,
\url{https://doi.org/10.1038/nrg.2016.49}. See also G.~A. Logsdon,
M.~R. Vollger and
E.~E. Eichler, ``Long-read human genome sequencing and its applications'',
\emph{Nature Reviews Genetics} 21 (2020), pp.~597--614,
\url{https://doi.org/10.1038/s41576-020-0236-x}.}

\subsection{Biological variation and measurement error}

Two homologous biological sequences can differ even if both are measured
perfectly. A single-nucleotide variant changes one nucleotide. An insertion or
deletion changes sequence length. Larger structural variants can duplicate,
invert, move or remove longer intervals. Recombination and horizontal transfer
can join segments with different histories. In transcript data, splicing
removes introns and joins exons, so a mature RNA or cDNA read can correspond to
multiple separated intervals in a genomic reference.

The algorithm sees many of these events through the same small vocabulary:
matches, substitutions and gaps. A mismatch may represent a true variant or a
base-calling error. A gap may represent a biological insertion or deletion, a
sequencing error, or a difference between the reference assembly and the
sample. A split alignment may represent a structural variant, a splice
junction, an assembly boundary or an experimental chimera. Interpretation
therefore requires metadata and context in addition to the aligned strings.
Figure~\ref{fig:variation-versus-error} traces how biological variation and
measurement error can produce the same observed disagreement.

\begin{figure}[ht]
  \centering
  \begin{tikzpicture}[
      scale=.76,
      transform shape,
      node distance=1.15cm and 1.35cm,
      >=Latex,
      box/.style={draw, rounded corners=2pt, align=center,
                  minimum height=8mm, text width=25mm, fill=black!3},
      cause/.style={draw, rounded corners=2pt, align=center,
                    minimum height=8mm, text width=27mm, fill=black!7}
    ]
    \node[box] (reference) {reference\\assembly};
    \node[box, below=of reference] (sample) {sample's true\\molecule};
    \node[cause, right=of sample] (variation) {biological\\variation};
    \node[cause, right=of variation] (sequencing) {sequencing and\\base calling};
    \node[box, right=of sequencing] (read) {observed\\query read};
    \node[box, above=of read] (comparison) {sequence\\comparison};

    \draw[->] (sample) -- node[above, font=\scriptsize] {differs from reference}
      (variation);
    \draw[->] (variation) -- node[above, font=\scriptsize] {true sample string}
      (sequencing);
    \draw[->] (sequencing) -- node[above, font=\scriptsize] {may add error}
      (read);
    \draw[->] (read) -- (comparison);
    \draw[->] (reference) -- (comparison);
  \end{tikzpicture}
  \caption{Two sources of disagreement meet at the same comparison. Biological
    variation separates the sample molecule from the reference, while
    sequencing and base calling can separate the observed read from that
    molecule. Alignment locates and describes disagreement. Metadata and
    experimental context are needed to interpret its cause.}
  \label{fig:variation-versus-error}
\end{figure}

Figure~\ref{fig:variation-versus-error} explains why exact string equality is
too strict and why alignment alone is not a variant caller. The comparison
stage must tolerate plausible differences. A later stage must then combine
the alignment with coverage, base qualities, sample relationships and other
evidence before assigning a biological interpretation.
Table~\ref{tab:alignment-ambiguity} maps common alignment patterns to the
distinct causes that remain unresolved at this stage.

\begin{table}[ht]
  \centering
  \small
  \begin{tabular}{p{0.16\linewidth}p{0.29\linewidth}p{0.29\linewidth}p{0.16\linewidth}}
    \toprule
    Observed pattern & Possible biological explanation &
    Possible observational or representation explanation &
    Alignment-level representation \\
    \midrule
    One symbol differs &
    single-nucleotide variant or divergence between homologues &
    base-calling error, consensus error or unresolved ambiguity &
    mismatch or scored substitution \\
    Extra or missing symbols &
    insertion or deletion in the sample lineage &
    insertion/deletion sequencing error or assembly gap &
    one or more gap operations \\
    Long discontinuity &
    structural variant, rearrangement or splice junction &
    chimeric read, reference discontinuity or assembly error &
    long gap or split alignment \\
    Several plausible locations &
    repeated or duplicated biological sequence &
    incomplete reference or query too short to disambiguate &
    competing mappings with similar scores \\
    \bottomrule
  \end{tabular}
  \caption{The same alignment pattern can have several causes. The table is a
    reminder that similarity search generates evidence. It does not by itself
    determine the biological event.}
  \label{tab:alignment-ambiguity}
\end{table}

\section{Reference and query are computational roles}
\label{sec:reference-query-roles}

Throughout this thesis, the \emph{reference} or \emph{target} is the sequence
collection that is prepared for search, often by building an index. The
\emph{query} is the sequence presented for comparison against that collection.
These are roles in an operation, not fixed classes of molecule. A chromosome
assembly can be a reference when reads are mapped to it and a query when it is
compared with another assembly. A read set can be searched against a reference
genome, against another read set to find overlaps, or against a database of
known sequences. BioSeqDB's SQL-facing similarity predicate follows the same
operational distinction by accepting a target sequence value and a query
sequence value.

The role assignment changes the engineering cost. Work spent once on a stable
reference index can be amortized over many queries. If both inputs change for
every comparison, as in an all-versus-all overlap workload, index construction
and reuse have a different balance. A database makes this reuse question
explicit because sequence values, indexes and query plans have managed
lifetimes rather than being unrelated files produced by a script.

Several coordinate conventions follow from the reference/query distinction.
Positions are usually reported in the coordinate system of the stored target.
The query can align in forward or reverse-complement orientation. Local
alignment can consume only a portion of either input. Gaps advance one sequence
without advancing the other. Split mappings can produce several target
intervals for one query. These facts explain why an alignment result is richer
than a Boolean answer and why downstream columns such as start, end,
orientation, score and edit description are useful.

\section{Why biological sequence comparison requires alignment}
\label{sec:why-alignment}

Exact equality answers whether two symbol strings are identical. Most
biological questions ask something different: whether the strings can be
explained as related sequences after a plausible number and arrangement of
changes. A useful comparison must therefore answer three connected questions:

\begin{enumerate}
  \item \emph{Where} can the query be placed relative to the target?
  \item \emph{How} do their symbols correspond, including substitutions and
    gaps?
  \item \emph{How credible} is that correspondence compared with alternatives?
\end{enumerate}

An \emph{alignment} makes the correspondence explicit by placing symbols in
columns and introducing gaps where one sequence advances without the other.
A scoring model rewards compatible symbol pairs and penalizes substitutions
and gaps. Global alignment attempts to explain both inputs end to end. Local
alignment searches for high-scoring similar regions. Semi-global and
free-end-gap variants model cases in which one input is expected to be
contained within another. Spliced and split alignment add structure for
discontinuous relationships.

The fully general dynamic-programming solution examines a matrix whose axes are
the two sequences. That formulation is precise, but its work grows with the
product of their lengths. Genome-scale tools therefore avoid evaluating most
matrix cells. Modern mappers commonly identify short exact or near-exact
anchors, connect mutually consistent anchors into chains, and perform detailed
alignment only around promising chains. Minimap2 is a prominent example of
this organization around seeding, chaining, and alignment.\footnote{H. Li,
``Minimap2: pairwise
alignment for nucleotide sequences'', \emph{Bioinformatics} 34(18), 2018,
pp.~3094--3100, \url{https://doi.org/10.1093/bioinformatics/bty191}.}
BioSeqDB's current alignment pipeline is minimap2-based, while replacing or
extending individual stages with CPU SIMD and GPU-oriented implementations.
The following chapter introduces these algorithmic families without assuming
that every method uses the same seed representation or the same final
alignment kernel.

The biological discussion above determines what those algorithms must
preserve. They must handle strand orientation, tolerate both substitutions and
length-changing differences, avoid being dominated by repeated sequence, and
retain coordinates that can be joined with annotations. They must also avoid
assigning biological meaning that the sequence comparison alone cannot
support.

\section{Why place sequence analysis inside a database?}
\label{sec:why-database-resident}

Sequence analysis is often presented as a file-to-file transformation. A FASTA
or FASTQ file enters a command-line program and an alignment file leaves it.
That model is effective for a single tool invocation, but a real study also
contains sample identifiers, collection conditions, organism labels, cohort
membership, reference versions, quality-control decisions and derived
annotations. When each stage writes another loosely coupled file, the
relationship between a sequence and its provenance is maintained by file
names, scripts and conventions outside the data itself.

A database-resident design changes the unit of work. Instead of asking an
aligner to process every record in a file, a relational query can first select
the relevant samples, join them with the chosen reference or query set, invoke
similarity search on the resulting sequence pairs, and aggregate or filter the
results. A schematic workload is:
\[
  \text{metadata selection}
  \;\longrightarrow\;
  \text{candidate sequence pairs}
  \;\longrightarrow\;
  \text{alignment}
  \;\longrightarrow\;
  \text{relational result}.
\]
This organization is valuable for five reasons.

\begin{enumerate}
  \item \textbf{Metadata and sequence remain addressable together.}
    Alignment results can be related to sample, experiment and reference
    records without reconstructing those relationships from external file
    naming schemes.
  \item \textbf{Cheap predicates can reduce expensive work.}
    Filtering by sample attributes or indexed identifiers before alignment can
    reduce the number of candidate comparisons. This is a planning
    opportunity, not a guarantee that every SQL formulation will be optimal.
  \item \textbf{Indexes acquire managed lifetimes.}
    A sequence index can be associated with stored data and reused across
    queries. Its construction, invalidation and selection become part of the
    database system rather than an implicit side effect of a script.
  \item \textbf{Results compose with ordinary relational operations.}
    Coordinates, scores and edit summaries can feed joins, grouping and
    aggregation. A biological workflow need not cross a process boundary after
    every comparison step.
  \item \textbf{Provenance and consistency have a common home.}
    Transactions and typed columns can keep a sequence value, its metadata and
    derived state synchronized. External resources and large data transfers
    still exist, but their relationship to stored records can be made
    explicit.
\end{enumerate}

BioSeqDB implements this idea as a PostgreSQL extension rather than as a
standalone aligner wrapped by SQL text. Current source exposes alphabet-aware
sequence values and an \texttt{mmsearch} similarity predicate. Integration
tests exercise searches over stored columns, constants, slices and joins.
PostgreSQL custom scans then provide execution paths that can batch sequence
work and return relational tuples. Later chapters explain the planner
rewrites, storage layout, index persistence, caching and vectorized aggregates
that make this interface more than a scalar function call.

Database residence does not imply that all computation is free or that no data
ever moves. Alignment remains expensive. GPU execution can require transfer to
device memory. External sequence providers still require network and I/O
boundaries. Large intermediate results can still dominate a query. The useful
claim is narrower: sequence storage, candidate selection, alignment execution
and downstream relational processing can participate in one managed system.
That integration creates opportunities for reuse, batching and provenance that
are difficult to obtain from a collection of unrelated file transformations.

\section{Transition from biology to algorithms}
\label{sec:biology-transition}

The biological problem can now be stated in computational terms. The input is
not simply ``DNA''. It is an alphabet-tagged reference or target, an
alphabet-compatible query, orientation and coordinate conventions, and
metadata describing the origin of both values. Differences can arise from
evolution, within-population variation, transcript structure, sequencing error
or incomplete representation. The required output is therefore not only a
similarity decision. It is a defensible correspondence between positions,
accompanied by enough score and coordinate information for later
interpretation.

The next chapter develops the algorithms that approximate this correspondence
at scale. It begins with exact dynamic programming to establish what an
alignment optimizes. It then explains seed-and-extend methods, minimizer-based
seeding and graph-based alternatives such as de Bruijn graphs. That progression
connects the biological sources of disagreement described here to the design
of BioSeqDB's processing pipeline: minimizer extraction proposes candidate
anchors, chaining assembles them into coherent placements, prefilters reject
unlikely pairs, and the final alignment step produces a base-level account of
matches, substitutions and gaps.
